\documentclass[12pt]{article}

\usepackage{sbc-template}
\usepackage{graphicx,url}
\usepackage[utf8]{inputenc}
\usepackage[T1]{fontenc}
\usepackage[brazilian]{babel}
\usepackage{geometry}
\usepackage{setspace}
\usepackage{titlesec}
\usepackage{indentfirst}
\sloppy
\title{Uma Análise Comparativa entre NeRFs Few-Shot e Per-Scene aplicadas em Conjuntos de Dados de Vistas Esparsas}
\author{Felipe G. da Silva\inst{1}}
\address{Departamento de Ciências da Computação e Estatística -- UNESP\\
  CEP: 15054-000 -- São José do Rio Preto -- SP -- Brazil
  \email{felipe-gomes.silva@unesp.br}
}

\begin{document} 
\maketitle
\begin{abstract}
This work compares hybrid Generalizable NeRFs and Fast Optmization Per-Scene trained models applied to sparse-view datasets, evaluating the impact of these feature extraction architectures on novel view synthesis performance. Through a quantitative analysis using visual fidelity (PSNR, SSIM, LPIPS) and computational efficiency metrics (convergence time, VRAM footprint, inference time), this study highlights the trade-offs between the global generalization of Transformers and the local processing efficiency of convolutions. Ultimately, it aims to guide the selection of 3D rendering models for practical scenarios often constrained by hardware limitations and data scarcity.
\end{abstract}

\begin{resumo}
  Esse trabalho apresenta uma comparação entre modelos generalizáveis via extratores de características e modelos de otimização rápida por cena aplicados em conjuntos de dados com vistas esparsas e avalia o impacto dessas arquiteturas de extração de características no desempenho em síntese de novas vistas. Por meio de uma análise quantitativa baseada em métricas de fidelidade visual (PSNR, SSIM e LPIPS) e de eficiência computacional (tempo de convergência, consumo de VRAM e tempo de inferência), o estudo busca evidenciar os trade-offs entre a capacidade de generalização global dos Transformers e a eficiência de processamento local das convoluções,  visando auxiliar na seleção de modelos de renderização 3D em cenários práticos, os quais são frequentemente marcados por restrições de hardware e escassez de dados.
\end{resumo}

\section{Introdução}
\label{sec:intro}

Modelos de campos de radiância neural (NeRF) consolidaram-se, nos anos recentes, como o estado da arte, tanto na indústria, quanto na academia, sobre geração de novas vistas \cite{kerbl3Dgaussians, Tancik2023, substance3dviewer}. Entretanto, as mais avançadas técnicas de renderização neural exigem uma grande quantidade de: (i) poder computacional; (ii) múltiplas vistas no conjunto de dados. Ambos os desafios tendem a impedir que aplicações mais simples, como escaneamento \textit{mobile} de objetos, provador virtual, geração de assets 3D e síntese em datasets pequenos, sejam aproveitados. É com tal problemática que projetos como Instant-NGP (\textit{Instant Neural Graphics Primitives}), MeRF (\textit{Memory-Efficient Neural Radiance Field}), Pixel-Nerf, Gnt (\textit{Generalizable Nerf Transformer}) \cite{muller2022, reiser2023merf, yu2021pixelnerfneuralradiancefields, t2023attentionnerfneeds}, entre outros, tentam expandir o estado da arte e, consequentemente, abordam temas análogos, como: (i) \textit{encoding} espacial de parâmetros \cite{muller2022, reiser2023merf}; (ii) uso de arquiteturas híbridas entre Convolutional Neural Networks (CNN) e Vision Transformers (ViT) \cite{schwarz2021grafgenerativeradiancefields, wu2021cvt, t2023attentionnerfneeds}; e (iii) pré-computação (\textit{baking}) dos modelos em estruturas de menor pegada de memória \cite{reiser2023merf, hedman2021bakingneuralradiancefields}.

Este trabalho tem, portanto, o objetivo de explorar o uso dos diferentes modelos de campos de radiância neural e suas respectivas técnicas de otimização, ao aplicá-las em alguns conjuntos de dados com vistas esparsas, promovendo uma análise exploratória entre desempenho computacional e qualidade dos volumes gerados. Para avaliar a fidelidade das reconstruções, serão utilizadas métricas de qualidade visual amplamente consagradas na literatura, como PSNR (\textit{Peak Signal-to-Noise Ratio}), SSIM (\textit{Structural Similarity Index Measure}) \cite{wang2004image} e LPIPS (\textit{Learned Perceptual Image Patch Similarity}) \cite{zhang2018unreasonable}. Paralelamente, a quantificação do custo operacional basear-se-á em métricas de desempenho sistêmico, incluindo o tempo de convergência no treinamento, a velocidade de inferência e a pegada de memória (\textit{memory footprint}), sobretudo o consumo de VRAM (\textit{Video Random Access Memory}). Embora trabalhos como \cite{niemeyer2022regnerf} e \cite{kim2023poserefinement} tenham avaliado NeRFs em regimes de dados limitados, a comparação sistemática entre modelos generalizáveis com extratores convolucionais versus baseados em atenção, sob condições controladas de hardware e splits reprodutíveis, permanece pouco explorada na literatura. Em suma, esta investigação fornece um panorama consolidado sobre os \textit{trade-offs} entre eficiência computacional e fidelidade visual, estabelecendo diretrizes claras para a viabilização e aplicação prática de modelos de síntese de novas vistas em cenários com restrições das mais diversas, especialmente de dados e de hardware.

\section{Fundamentação Teórica}
\label{sec:fund}

\teorical{Conjuntos de Dados de Vistas Esparsas} Essa nomenclatura refere-se, neste trabalho, aos conjuntos de dados imagéticos que possuem pouca (ou nenhuma) redundância e/ou variabilidade de diferentes poses de câmera referentes ao mesmo objeto. Áreas como histologia, holotomografia e microscopia confocal de fluorescência ilustram bem esse problema: a aquisição de imagens multi-vista é fisicamente limitada, biologicamente invasiva e estruturalmente esparsa por natureza \cite{XUE2021101816}.


\teorical{Campos de Radiância Neural} Campos de Radiância Neural (\textit{Neural Radiance Fields}, NeRF) utilizam-se de redes neurais do tipo MLP (\textit{Multilayer Perceptron}) para aproximar uma função contínua no espaço 5D. 
Essa função recebe como entrada uma posição espacial $x = (x, y, z)$ e uma direção de visualização $d = (\theta, \phi)$, retornando a densidade volumétrica $\sigma(x)$ e a cor associada $c = (r, g, b)$:
\[
F_\theta(x, d) = F_\theta\prime \circ \gamma(x, d) \rightarrow (\sigma, c)
\]
onde $\theta$ representa os pesos da rede neural e $\gamma$ um codificador posicional.

\teorical{Representações de Cena} Trabalhos subsequentes introduziram representações híbridas que desacoplam a estrutura espacial da rede neural, dado a viabilidade de alterar-se o \textit{positional encoding} via codificação de fourier por um sistema de \textit{hash encoding}. O Instant-NGP \cite{muller2022} propõe uma codificação por tabela de hash multiescala, que mapeia coordenadas espaciais em vetores de características de forma eficiente, reduzindo drasticamente o tempo de otimização por cena. O MERF \cite{reiser2023merf} estende essa ideia com representações tridimenseionais esparsas orientadas à renderização em tempo real.

\teorical{Extrator de Características} A capacidade de generalizar para novas cenas a partir de poucas vistas depende criticamente do extrator de características, responsável por codificar as imagens de entrada em representações que a rede neural possa condicionar. O PixelNeRF \cite{yu2021pixelnerfneuralradiancefields} utiliza uma ResNet pré-treinada para extrair mapas de características por imagem; dado um ponto 3D, sua projeção nas vistas de entrada recupera os vetores de características espaciais correspondentes, que são agregados e passados à MLP de radiância. Essa abordagem é local e opera por pixel, sem modelar explicitamente relações entre vistas. O GNT \cite{t2023attentionnerfneeds} substitui a MLP de radiância por um transformador que processa amostras ao longo do raio e agrega informação entre vistas por meio de atenção cruzada, eliminando a dependência de uma função de volume explícita.

\teorical{Renderiznado o Volume} Seja um campo de radiância neural $F_\theta$, a cor esperada de um pixel é computada através de renderização volumétrica. Dado um raio $\mathbf{r}(t) = \mathbf{o} + t\mathbf{d}$ com limites de profundidade $(t_n, t_f)$, a cor resultante é dada pela integral:

\[
\hat{C}(\mathbf{r}) = \int_{t_n}^{t_f} T(t) \sigma(\mathbf{r}(t)) \mathbf{c}(\mathbf{r}(t), \mathbf{d}) \, dt
\]
onde $T(t) = \exp\left(-\int_{t_n}^t \sigma(\mathbf{r}(s)) \, ds\right)$ denota a transmitância acumulada ao longo do raio, ou seja, a probabilidade de o raio não atingir nenhuma partícula obstrutiva até a profundidade $t$. Na prática, por ser uma função contínua, essa integral é aproximada numericamente via \textit{alpha compositing} sobre amostras discretas interpoladas ao longo de cada raio durante o treinamento.

\teorical{NeRFs Few-Shot vs. Per-Scene} Com o objetivo de reduzir a necessidade de múltiplas vistas durante o treinamento, surgiram abordagens conhecidas como NeRFs \textit{few-shot}, capazes de generalizar a cena e produzir novas vistas a partir de poucas imagens, utilizando-se de diversos conceitos distintos aplicados aos componentes NeRF supracitados. Estas abordagens contrastam frontalmente com os métodos \textit{per-scene}. Enquanto modelos \textit{per-scene} (como Instant-NGP) otimizam uma estrutura de dados de forma isolada e colapsam sob a escassez de vistas — gerando graves artefatos visuais e ambiguidades geométricas —, arquiteturas \textit{few-shot} e generalizáveis incorporam \textit{priors} estruturais e semânticos aprendidos de grandes bases de dados. Essa distinção arquitetural torna-se indispensável em cenários onde a extração é severamente restrita. Em domínios de imagens patológicas e histologia, por exemplo, a capacidade de inferir uma geometria coerente a partir de amostragem esparsa é fundamental, viabilizando o uso desses modelos generativos não apenas para síntese, mas como mecanismos robustos de aumento de dados (data augmentation) para mitigar desbalanceamentos em pipelines complexos de classificação médica.

\section{Metodologia} 
\label{sec:metodology}

% The first page must display the paper title, the name and address of the
% authors, the abstract in English and ``resumo'' in Portuguese (``resumos'' are
% required only for papers written in Portuguese). The title must be centered
% over the whole page, in 16 point boldface font and with 12 points of space
  % before itself. Author names must be centered in 12 point font, bold, all of
% them disposed in the same line, separated by commas and with 12 points of
% space after the title. Addresses must be centered in 12 point font, also with
% 12 points of space after the authors' names. E-mail addresses should be
% written using font Courier New, 10 point nominal size, with 6 points of space
% before and 6 points of space after.

% The abstract and ``resumo'' (if is the case) must be in 12 point Times font,
% indented 0.8cm on both sides. The word \textbf{Abstract} and \textbf{Resumo},
% should be written in boldface and must precede the text.

Esta seção contempla a metodologia adotada para a realização dos experimentos, incluindo a descrição dos conjuntos de dados utilizados, os modelos de NeRF selecionados para o estudo comparativo, as métricas de avaliação aplicadas e o ambiente computacional onde os testes foram conduzidos. 

Para viabilizar a comparação entre os modelos, implementou-se, individualmente, (caso não houvesse presença prévia) a integração do modelo proposto com a ferramenta NeRFStudio. A seguir, apresentam-se os modelos selecionados.

\subsection{Modelos Selecionados}
Neste trabalho, são avaliados modelos representativos para reconstrução de campos de radiância neural em cenários de poucas vistas (\textit{few-shot}) e generalizáveis, contemplando diferentes paradigmas arquiteturais e estratégias de aprendizado. Especificamente, consideram-se: 

\begin{enumerate}
\item PixelNeRF \cite{yu2021pixelnerfneuralradiancefields}, que utiliza redes neurais convolucionais para extrair características espaciais das imagens e inseri-las no renderizador através de interpolação e projeção; 

\item Generalizable NeRF Transformer (GNT) \cite{t2023attentionnerfneeds}, que emprega mecanismos de atenção (\textit{Transformers}) em ambos os componentes NeRF, substituindo tanto o extrator de características quanto o renderizador diferenciável por módulos baseados em atenção cruzada; 
\end{enumerate}

Para fins de linha de base (baseline) comparativa, este estudo também avalia métodos \textit{per-scene} sob regime de dados limitados: 

\begin{enumerate}
\item Instant-NGP \cite{muller2022}, que substitui a MLP tradicional por uma codificação em tabela de hash multiescala, permitindo otimização por cena em tempo significativamente reduzido; e

\item MERF \cite{reiser2023merf}, que estende representações baseadas em hash com estruturas esparsas orientadas à renderização em tempo real.
\end{enumerate}

A seleção desses modelos visa proporcionar uma análise abrangente entre diferentes abordagens, cobrindo métodos baseados em CNNs, transformadores, reconstrução volumétrica explícita e modelos generativos baseados em representações por hash \textit{per-scene}. Essa diversidade permite avaliar não apenas a qualidade de reconstrução em cenários com vistas limitadas, mas também o compromisso entre generalização entre cenas e custo de otimização, bem como o impacto de estratégias distintas de extração de características nos três regimes de avaliação considerados, como pode ser visto na seção \ref{sec:stdin}.

Pode-se perguntar o por quê de não utilizar-se outros tipos de modelos. A seguir, destaca-se motivações específicas para a não seleção de grandes modelos possíveis:
\begin{itemize}
\item RegNeRF: o próprio trabalho % TODO
\item MVSNeRF
\item MIPNeRF-360

\end{itemize}


\subsection{Implementação no Nerfstudio}
Os modelos foram implementados no framework \textit{Nerfstudio} \cite{Tancik2023}, seguindo sua arquitetura modular. Para cada abordagem, foram definidos componentes específicos de \textit{DataManager}, \textit{Model}, \textit{Pipeline}, \textit{Trainer} e configurações associadas, aproveitando ao máximo as abstrações pré-existentes do framework e realizando adaptações pontuais para garantir compatibilidade com os modelos selecionados. A adoção de um framework unificado assegura equivalência no protocolo de avaliação entre os modelos, eliminando variáveis de implementação que poderiam comprometer a comparação, visto que o \textit{Nerfstudio} oferece, ademais, um conjunto de utilitários aproveitados neste trabalho:

\begin{enumerate}
\item visualização em tempo real por meio de interface gráfica interativa, permitindo monitorar o treinamento e inspecionar qualitativamente os resultados de renderização durante o processo;

\item aceleração via GPU com suporte a CUDA, integrada à biblioteca \textit{tiny-cuda-nn} \cite{tinycudann}, que fornece implementações otimizadas de MLPs e codificações de entrada diretamente em kernels CUDA, com ganho de desempenho reportado entre 5$\times$ e 30$\times$ em relação a implementações PyTorch convencionais, dependendo da arquitetura e do cenário de uso; e

\item infraestrutura de logging e exportação de métricas, incluindo salvamento de checkpoints e integração com \textit{Weights \& Biases}, que padroniza o registro dos experimentos entre os diferentes modelos.
\end{enumerate}

Dentre os modelos avaliados, o Instant-NGP e o MERF encontram-se nativamente integrados ao \textit{Nerfstudio}, dispensando adaptações de implementação. Os demais modelos --- PixelNeRF e GNT  --- foram portados para o framework a partir de suas implementações originais, com adaptações nos componentes de \textit{DataManager} e \textit{Pipeline} para conformidade com a interface do \textit{Nerfstudio}. No caso do MERF, adicionalmente, foi utilizada a extensão \textit{MERFStudio} \cite{song2024cityonwebrealtimeneuralrendering}, disponível em \url{https://github.com/USTC3DV/MERFStudio}.

As implementações foram guiadas pelas descrições metodológicas dos trabalhos originais, mantendo equivalência funcional sempre que possível. A implementação completa encontra-se disponível em \url{https://github.com/felipe-gsilva/few-shot-nerfs}.

% TODO talvez explicitar mais o código

\subsection{Protocolo Experimental}
\label{sec:stdin}

Os experimentos foram conduzidos sobre o dataset LLFF (\textit{Local Light Field Fusion}) \cite{mildenhall2019locallightfieldfusion}, composto por cenas de objetos reais capturadas com câmera em movimento frontal. Seguindo o protocolo estabelecido por RegNeRF \cite{niemeyer2022regnerf}, reservam-se para teste as imagens cujo índice é múltiplo de 8 (llffhold=8), totalizando entre 5 e 8 imagens de teste por cena. As N vistas de treinamento ($N \in {3, 6, 10}$) são amostradas uniformemente das imagens remanescentes.

A seleção das vistas de treino foi realizada por duas estratégias distintas: (i) \textit{uniform}, em que as vistas são amostradas por interpolação linear sobre os índices do dataset (\texttt{np.linspace}), distribuindo-as uniformemente ao longo da trajetória da câmera; e (ii) \textit{random}, em que as vistas são sorteadas aleatoriamente com semente fixa (\texttt{seed = 42}), garantindo reprodutibilidade. A comparação entre as duas estratégias permite avaliar a sensibilidade dos modelos à distribuição espacial das vistas de entrada.

As imagens restantes foram divididas proporcionalmente entre conjuntos de validação e teste, também com semente fixa, assegurando que as vistas de teste sejam idênticas entre todos os modelos e estratégias de split, eliminando variáveis de avaliação que poderiam comprometer a comparação. As métricas finais são reportadas exclusivamente sobre o conjunto de teste, sem sobreposição com as vistas de treino. Todos os experimentos foram executados em uma GPU NVIDIA GeForce RTX 4060 Ti (16GB de VRAM), com CUDA 12.4 e driver 580.126.09.
Para garantir reprodutibilidade em diferentes arquiteturas, optou-se por utilizar a ferramenta docker, a qual permite instanciar a GPU virtualmente e viabilizar o uso do torch, tiny-cuda-nn e cuda 12.4 entre diferentes arquiteturas de GPU.


\subsection{Métricas de Avaliação}
\label{sec:metricas}

A avaliação dos modelos contempla três dimensões complementares: qualidade visual, eficiência computacional e pegada de memória.

Para a qualidade visual, são empregadas três métricas amplamente adotadas na literatura de síntese de novas vistas: (i) PSNR (\textit{Peak Signal-to-Noise Ratio}), que mede a fidelidade pixel a pixel entre a imagem renderizada e a imagem de referência em escala logarítmica, sendo sensível a erros de brilho e cor; (ii) SSIM (\textit{Structural Similarity Index Measure}) \cite{wang2004image}, que avalia a similaridade estrutural entre imagens considerando luminância, contraste e estrutura, sendo mais correlacionado com a percepção humana do que o PSNR; e (iii) LPIPS (\textit{Learned Perceptual Image Patch Similarity}) \cite{zhang2018unreasonable} usando a rede vgg (Visual Geometry Group) como extrator de características \cite{niemeyer2022regnerf} para quantificar a distância perceptual entre imagens (valores menores indicam maior similaridade perceptual). Para a eficiência computacional, é medida a taxa de renderização em quadros por segundo (\textit{fps}), avaliada durante a inferência sobre as câmeras do conjunto de teste. Para a pegada de memória, é reportado o footprint de parâmetros do modelo, calculado como a soma do produto entre o número de parâmetros e o tamanho de cada elemento em bytes. Também comparou-se os respectivos tamanhos de parâmetros com a quantidade de memória alocada no torch naquele dado momento. Como o ambiente de GPU é controlado (reservado pelo ambiente docker), possíveis outros processos simultâneos não adicionaram ruído às métricas de avaliação.


As avaliações de qualidade visual são realizadas sobre o conjunto de teste gerado pelo protocolo de split descrito na Seção \ref{sec:stdin}, utilizando o \textit{dataloader} de avaliação do \textit{Nerfstudio} (\texttt{fixed\_indices\_eval\_dataloader}), que itera sobre câmeras fixas e compara as imagens renderizadas (\texttt{outputs["rgb"]}) com as imagens de referência (\texttt{batch["image"]}) sem amostras das vistas de treino.

\section{Experimentos e Resultados}

\subsection{Análise Quantitativa}

\subsection{Análise Qualitativa}

\subsection{Trade-offs Computacionais}

% In some conferences, the papers are published on CD-ROM while only the
% abstract is published in the printed Proceedings. In this case, authors are
% invited to prepare two final versions of the paper. One, complete, to be
% published on the CD and the other, containing only the first page, with
% abstract and ``resumo'' (for papers in Portuguese).

% TODO

\section{Conclusões}

% Section titles must be in boldface, 13pt, flush left. There should be an extra
% 12 pt of space before each title. Section numbering is optional. The first
% paragraph of each section should not be indented, while the first lines of
% subsequent paragraphs should be indented by 1.27 cm.

\bibliographystyle{sbc}
\bibliography{sbc-template}

\end{document}
